% TOP_PROBLEM: {"id": 4600019, "title": "Invariant measures and subsystems", "category_id": 11, "category": {"id": 11, "name": "dynamical_systems", "display_name": "Dynamical Systems", "description": "Problems about long-term behavior of deterministic systems, Hamiltonian dynamics, and periodic orbits.", "slug": "dynamical-systems", "order_index": 11, "created_at": "2026-07-31T15:26:25.671Z"}, "status": "open", "problem_number": "AMR-045-0019", "set_id": 13, "statusQualification": "Restores the explicit three-dot system X in source Problem 14.3; no Haar, ergodicity, or full-support restriction is imposed."}
% FIRST_POSTED: 2026-10-01
% ENTRY_KIND: statement_only
% UnsolvedMath ID 4600019; source catalogue code AMR-045-0019
% !TeX program = lualatex
% Definition and sources only; this file is not an LLM solution attempt.
\documentclass[11pt,a4paper]{article}
\usepackage[margin=25mm]{geometry}
\usepackage{fontspec}
\setmainfont{lmroman10-regular.otf}[BoldFont=lmroman10-bold.otf,
 ItalicFont=lmroman10-italic.otf,BoldItalicFont=lmroman10-bolditalic.otf]
\usepackage{amsmath,amssymb,mathtools}
\usepackage{unicode-math}
\setmathfont{latinmodern-math.otf}
\AtBeginDocument{\let\setminus\smallsetminus}
\usepackage{xurl,hyperref}
\hypersetup{colorlinks=true,urlcolor=blue}
\setcounter{secnumdepth}{0}
\setlength{\parindent}{0pt}
\setlength{\parskip}{5pt}
\setlength{\emergencystretch}{3em}
\begin{document}
\section*{Invariant measures and subsystems}\label{4600019}
{\small UnsolvedMath ID 4600019 · AMR-045-0019 · Dynamical Systems · AMR Open Problem Lists}
\subsection{Definitions and mathematical statement}
Consider the compact binary configuration space
\[
 X=\{x\in(\mathbb Z/2\mathbb Z)^{\mathbb Z^2}:
       x_{i,j}+x_{i+1,j}+x_{i,j+1}=0\pmod2
       \text{ for every }(i,j)\}.
\]
For $v\in\mathbb Z^2$, the shift $\alpha^v$ translates coordinates by
$v$. The problem has two classification targets. First, describe all
Borel probabilities $\mu$ on $X$ satisfying
$\mu((\alpha^v)^{-1}E)=\mu(E)$ for every Borel $E$ and every $v$.
Second, describe all nonempty closed sets $Y\subset X$ with
$\alpha^vY=Y$ for every $v\in\mathbb Z^2$.

Such a $Y$ is called a subsystem, with the restricted shift action.
It need not be a subgroup of the compact group $X$, and it need not
have a finite forbidden-pattern presentation. Likewise, the probabilities
need not be Haar measures on invariant subgroups, need not be ergodic,
and need not have full support. A measure classification should in
particular identify the ergodic possibilities; general invariant measures
then include their ergodic decompositions.

The defining parity rule fixes the ambient system throughout, so the
question asks for its invariant probability simplex and its closed
invariant subsets, rather than for arbitrary binary planar shifts.
Finite orbits and proper subsystems are included. Source Problem 14.3
concerns this Ledrappier three-dot system, despite the catalogue's broad
reference to the Furstenberg discussion earlier in that section.
\subsection{Short English statement}
Classify all invariant probabilities and closed invariant subsets of the Ledrappier three-dot shift.
\subsection{Sources}
\begin{itemize}
\item[S1] ULAM.AI and the UnsolvedMath Contributors, supplied problems.json, record 4600019 (AMR-045-0019), AMR Open Problem Lists. Catalogue identity and source transcription; CC BY 4.0.
\url{https://huggingface.co/datasets/ulamai/UnsolvedMath}.
\item[S2] Boyle - Open Problems in Symbolic Dynamics (2008), Problem/Question/Conjecture 14.3. Original source indexed by AMR-045-0019.
\url{https://www.math.umd.edu/~mboyle/papers/openfinalsub3nov2007.pdf}.
\end{itemize}
\end{document}
